% page 206 of the facsimile (image p-205 of the scan)
% block 160.0 x 246.3 mm ; left margin 8.0 mm
\pgc{-12.700mm}{0.000mm}{
\l{\cel{3}198.}
\l{}
\l{\sou{goeleto}. (\sou{nav.})\cel{2}Navo lejera, du-masta. - DFIS.}
\l{}
\l{\sou{+golfo.}\cel{2}Ludo quan partoprenas 2 o 4 personi, en qua li}
\l{\cel{3}uzas quaza biliardo-buli, mikra e tre harda, quin onu}
\l{\cel{3}duktas, per bastono-frapi, rulige sur sulo, ad 9 o}
\l{\cel{3}18 trui en ica. - DEFIRS.}
\l{}
\l{\sou{gonagro.(patol.})\cel{2}Goto di la genuo. - DEF.}
\l{}
\l{\sou{gondo}. Fer-peco axala, kudatra e pivotoza, qua sustenas la}
\l{\cel{3}bendo sur qua rotacas la pordo, e c. - F.}
\l{}
\l{\sou{gondolo}. I. Vehilo suspendita sub aerostato (direktebla o ne).}
\l{\cel{3}- II. Speco di batelo uzata en la laguni, precipue en la}
\l{\cel{3}urbo Venezia (Italia). - DEFIRS.}
\l{}
\l{\sou{gonfalono}. I. Banero militala, suspendita anla supra parto}
\l{\cel{3}di lanco, di standardo. - II. Banero ekleziala (katolika)}
\l{\cel{3}sub qua (dum la Mezepoko) su asemblis la vasali di la Ekle-}
\l{\cel{3}zio. - F.}
\l{}
\l{\sou{goniometrio}. Cienco pri la mezuro di la anguli. - DEF.}
\l{}
\l{\sou{goniometro}. (\sou{geom.}) Instrumento por mezurar la anguli (di la}
\l{\cel{3}kristali, precipue). - DEF.}
\l{}
\l{\sou{gonokocio}. (\sou{patol.}) Nomo di la ensemblo de la morbi quin pro-}
\l{\cel{3}duktas la gonokoko od en qui la gonokoko pleas rolo certa:}
\l{\cel{3}blenoragii, metriti, oftalmii blenoragiala, e c.}
\l{}
\l{\sou{gonoreo}. (\sou{patol}.) Ekfluajo de la membrano genitala ed urin-}
\l{\cel{3}ifala. - DEFIRS.}
\l{}
\l{\sou{gorgear.}\cel{2}(\sou{netrans.}) Bruisar quale aquo qua falas de tekto-}
\l{\cel{3}pluvkanalo. - DEFI.}
\l{}
\l{\sou{gorilo}. (\sou{zool.})\cel{2}Simio di Afrika, analoga a chimpanzeo.}
\l{\cel{3}- L.\cel{1}\sou{gorilla gina}. - DEFIRS.}
\l{}
\l{\sou{gorjereto}. (\sou{kirurg.}) Instrumento kanaloza, kun kanalo streta,}
\l{\cel{3}por litotomiar od operacar ulu pri fistulo.- DEFIRS.}
\l{}
\l{\sou{goto.(patol.}) Morbo qua atakas ordinare la artiki, iras}
\l{\cel{3}de ica ad ita, ed ulakaze a la viceri. - EFIS.}
\l{}
\l{\sou{gotika}. I. (\sou{arkitekt.}) Ogivala, qua sucedis la Romanala. -}
\l{\cel{3}II. -(\sou{pri literi}). Perpendikla e di qua la konturi esas}
\l{\cel{3}quaze facetoza ed ornita ye pinti, hoki, ankore uzata}
\l{\cel{3}en la libri Germana. - DEFIRS.}
\l{}
\l{\sou{grabar}. (\sou{trans}.) I. Trasar sur materio harda (per stilo, cizel-}
\l{\cel{3}ilo) (figuro, epigrafo). - II. Trasar (per stilo, cizelilo)}
\l{\cel{3}sur plako de metalo o de ligno, la kopiuro de pikturo, de}
\l{\cel{3}desegnuro, de muziko-texto, por reproduktar lu mult-exem-}
\l{\cel{3}plere. - DEFIRS.}
}
